\newpage
\begin{center}
\large\textbf{Минобрнауки России}

\large\textbf{Юго-Западный государственный университет}
\vskip 1em
\normalsize{Кафедра программной инженерии}
\vskip 1em
\ifВКР{
        \begin{flushright}
        \begin{tabular}{p{.4\textwidth}}
        \centrow УТВЕРЖДАЮ: \\
        \centrow Заведующий кафедрой \\
        \hrulefill \\
        \setarstrut{\footnotesize}
        \centrow\footnotesize{(подпись, инициалы, фамилия)}\\
        \restorearstrut
        «\underline{\hspace{1cm}}»
        \underline{\hspace{3cm}}
        20\underline{\hspace{1cm}} г.\\
        \end{tabular}
        \end{flushright}
        }\fi
\end{center}
\vspace{1em}
  \begin{center}
  \large
\ifВКР{
ЗАДАНИЕ НА ВЫПУСКНУЮ КВАЛИФИКАЦИОННУЮ РАБОТУ
  ПО ПРОГРАММЕ БАКАЛАВРИАТА}
  \else
ЗАДАНИЕ НА КУРСОВУЮ РАБОТУ (ПРОЕКТ)
\fi
\normalsize
  \end{center}
\vspace{1em}
{\parindent0pt
  Студента \АвторРод, шифр\ \Шифр, группа \Группа
  
1. Тема «\Тема\ \ТемаВтораяСтрока»
\ifВКР{
утверждена приказом ректора ЮЗГУ от \ДатаПриказа\ № \НомерПриказа
}\fi.

2. Срок предоставления работы к защите \СрокПредоставления

3. Исходные данные для создания программной системы:

3.1. Перечень решаемых задач:}

\renewcommand\labelenumi{\theenumi)}

\begin{enumerate}
\item проанализировать предметную область управления проектами и сопоставить функциональные возможности существующих ИС~УП, доступных российским пользователям;
\item сформировать техническое задание на разработку веб-ориентированной программно-информационной системы для управления проектами;
\item спроектировать архитектуру программной системы, схему базы данных и программный интерфейс взаимодействия клиентской и серверной частей;
\item сконструировать клиентскую и серверную части программной системы, реализовать диаграмму Ганта с поддержкой четырех типов зависимостей задач и разграничение прав участников;
\item провести системное тестирование программной системы.
\end{enumerate}

{\parindent0pt
  3.2. Входные данные и требуемые результаты для программы:}

\begin{enumerate}
\item Входными данными для программной системы являются: учетные данные пользователей (логин, адрес электронной почты, ФИО, пароль), параметры проектов (название, описание, даты начала и завершения), параметры задач (название, описание, даты начала и завершения), сведения о зависимостях между задачами (предшественник, преемник, тип связи), состав участников проектов с указанием ролей, настройки цветовой темы интерфейса.
\item Выходными данными для программной системы являются: реестр проектов пользователя, диаграмма Ганта проекта со списком задач, стрелками зависимостей, линией текущей даты и индикацией истекшей части полос, перечень участников проекта с их ролями, цветовая палитра интерфейса, сообщения о результатах операций над данными.
\end{enumerate}

{\parindent0pt

  4. Содержание работы (по разделам):

  4.1. Введение.

  4.2. Анализ предметной области: управление проектами как область деятельности, методологии управления проектами, обзор информационных систем для управления проектами, сравнительный анализ существующих решений.

4.3. Техническое задание: основание для разработки, цель и назначение разработки, требования к функциональности, требования к интерфейсу, моделирование вариантов использования, требования к надежности и информационной безопасности, требования к программной и технической совместимости, требования к оформлению документации.

4.4. Технический проект: общая характеристика проектного решения, выбор средств и технологий реализации, архитектурное построение системы, проект базы данных.

4.5. Рабочий проект: спецификация компонентов и классов программной системы, системное тестирование программной системы.

4.6. Заключение.

4.7. Список использованных источников.

5. Перечень графического материала:

\списокПлакатов

\vskip 2em
\begin{tabular}{p{6.8cm}C{3.8cm}C{4.8cm}}
Руководитель \ifВКР{ВКР}\else работы (проекта) \fi & \lhrulefill{\fill} & \fillcenter\Руководитель\\
\setarstrut{\footnotesize}
& \footnotesize{(подпись, дата)} & \footnotesize{(инициалы, фамилия)}\\
\restorearstrut
Задание принял к исполнению & \lhrulefill{\fill} & \fillcenter\Автор\\
\setarstrut{\footnotesize}
& \footnotesize{(подпись, дата)} & \footnotesize{(инициалы, фамилия)}\\
\restorearstrut
\end{tabular}
}

\renewcommand\labelenumi{\theenumi.}
